\section{Implementation and Evaluation}
\label{sec:evaluation}

\noindent\textbf{Builds and measurement.}
The current-build series is based on \texttt{f856c7b}: three
100-hop runs per continuation mode and three runs of the
representative paused-adaptation treatment P. All nine runs
share the same node build. Additional timing probes are confined
to \texttt{payctl} and are enabled in all six chain runs, without
changes to production node paths. The complete N/R/C/P matrix
on \texttt{721800c} remains archived; historical-reader measurements
retain \texttt{721800c}, and the injected-latency study retains
\texttt{9879f64}. Chain summaries use medians and ranges of three
run totals; mixed-load summaries use medians of three run-level
statistics.

\noindent\textbf{Setup and measurement.}
Experiments use one Mac Studio M4 Max with 16 CPU cores, 64\,GiB RAM,
macOS~26.6.2, and Go~1.27.1. The project overlay of CometBFT~v0.38.26
implements the identity and installation rules. Network trials use
synchronous bbolt member commits before signing; committee application
state, BlockStore, and gateways are disk-backed. Members use
\texttt{GOMAXPROCS=16} and \texttt{GOGC=200}; commit, flush, and gossip
settings are 250, 10, and 10\,ms. Network trials copy the same unrun genesis and keys, reset genesis time,
and use sufficient fixed funding without top-up.

Certified receipt runs from first actual send to recipient verification
and atomic local recording. Public and member observations include
block following and driver polling; the chain experiment separately records
application-commit completion. Unless stated otherwise, summaries take
the median of three run-level statistics, not pooled transaction
quantiles; ranges describe runs, not confidence intervals. The timing trials do not measure restart recovery, sustainable
throughput, or minimum capital. The persistent restart checks below
are separate functional regressions.


\subsection{Continuous Spending Without Per-Hop Settlement}
\label{sec:eval-continuation}

Nine services support three independent NoSync wallets rotating 100~CAL through online
100-hop chains, using synchronous member commits and user-funded
final FUEL. The three pairs run in fast/wait, wait/fast, fast/wait order after a uniform three-second warm-up. Fast mode proceeds after certified receipt; wait mode
also requires the driver's observation of wallet-verified public
execution between hops. Both measure
$\mathsf{Sent}_1 \rightarrow \mathsf{Ready}_{100}$, including
construction of hops 2--100 and retries, but excluding first-payment
construction. Wait mode adds 99 intermediate
waits, not the final payment's public confirmation.

Median chain times are 4.707\,s and 64.479\,s, respectively,
a 13.70-fold ratio of medians; run ranges are 4.700--4.726\,s
and 64.387--64.670\,s. Fast uses 300 attempts for 300 payments,
whereas wait uses 414. All 297 fast successors use the preceding
TXCer and are first sent before the earliest parent
application-commit completion across the four replicas, with a
minimum lead of 183.439\,ms. Actual child-before-source obligations
number 2, 4, and 1 in the three fast runs; early sending alone is
not counted as out-of-order execution. All 600 payments close and
pass the accounting, input, replica, fee, and outbox audits.
Fast-mode all-payment member closure has a median of 5.526\,s
from the first send. Last-payment public observation has a median of 5.477\,s. Each chain spends 9,400~FUEL: 8,400 in rewards and 1,000 burned. For the fixed
transaction shape, the TXCer is 779~B and requests are 1,979~B
initially and 2,766~B thereafter.

\noindent\textbf{Client timing.}
For target height $h$, \texttt{BlockData} obtains the signed header
at $h+1$ before fetching the block and results at $h$. The RPC can
supply the latest-height seen commit without waiting for $h+2$.
From application-commit completion at the fetch source,
\texttt{committee0}, commit-to-complete-fetch includes successor
evidence, replica progress, polling, and retrieval, not pure polling
or query overhead. Its wait-mode P50/P95 are 285.884/314.509\,ms.
Separate P95s for post-fetch verification, preparation/local recording,
and post-transaction callback-to-driver observation are 0.236, 1.343,
and 4.713\,ms. The callback does not await a durability flush.
Supplement~S4 specifies checkpoints, clocks, matching, and aggregation.
The chain ratio compares end-to-end modes, including retries and
evidence observation; negative fast observation-to-build differences
indicate overlap.

\begin{figure}[t]
  \centering
  \includegraphics[width=\columnwidth]{continuation-final.pdf}
  \caption{Current-build 100-hop chains on \texttt{f856c7b}, including
    client timing probes: three runs per mode, synchronous members and
    NoSync wallets. Dots are run totals; bars mark medians. Only wait
    mode requires public observation between hops.}
  \label{fig:continuation}
\end{figure}

\subsection{Source Resolution Under Representative Mixed Load}
\label{sec:eval-mixed}
\label{sec:eval-liabilities}
\label{sec:eval-repair}

The fourteen-service P treatment schedules 7,000 ordinary payments
at 100 payments/s over 70\,s, alongside eight independent
cross-organization parent--child pairs, using authorized
organization-funded FUEL. Four committee nodes and two organizations with four members and one gateway each give fourteen services. Ordinary wallets use NoSync; scenario wallets commit synchronously. The driver limits fast requests to 64 and outstanding payments to 256. Pairs start ten seconds after driver launch, one second apart, and scenario wallets follow blocks for at least 55\,s. Parent publication and \texttt{INSTALL}
are withheld, automatic source resubmission is disabled, and all
adaptation endpoints are paused; consensus and compensation
decisions remain active. Each late source is released after its
compensation is observed. Ordinary traffic spans the injected
fault and recovery intervals.

All 21,048 payments across the three runs close, including
21,000 ordinary payments and 24 parent--child pairs. Each run
pays and recovers 800~CAL. Before adaptation resumes, all four
committee replicas report all eight targets \texttt{Recovered},
with their representation \texttt{Committed} and
\texttt{Materialized} flags both false. These 96 replica--target
checks establish compensation and repayment before adaptation
resumption in the measured executions. Within each run,
committee state hashes agree at a common height, and CAL/FUEL,
input, fee, reservation, and outbox audits pass. Per-run fees are
659,544~FUEL: 589,384 in rewards and 70,160 burned.

Table~\ref{tab:current-p} reports ordinary receipt latency together
with pre-send backpressure and drain. Median run-level receipt
P50/P95 are 54.538/147.549\,ms; public and member-observation P95s
are 2.187 and 2.255\,s. All three runs reach the 256-outstanding
limit. Dispatch P95 reaches 780.727\,ms and scheduled-to-receipt
P95 reaches 851.742\,ms; the latter is computed per payment,
not by adding quantiles. Actual send spans are 70.080--70.771\,s,
followed by 0.955--1.115\,s of ordinary-payment drain. This
representative P series demonstrates economic closure during
ordinary traffic with adaptation unavailable; it is not a
current-build comparison of interference across N/R/C/P.
\begin{table*}[t]
\centering
\caption{Ordinary-payment measurements for representative P runs
on \texttt{f856c7b}. Each run contains 7,000 ordinary payments
plus eight parent--child pairs. Receipt starts at actual send;
dispatch spans scheduled to actual send; scheduled-to-receipt
is calculated directly per payment. Drain starts at the last
ordinary send. The final row takes the median of each column's
three run-level entries.}
\label{tab:current-p}
\label{tab:eval-mixed}
\small
\setlength{\tabcolsep}{5pt}
\begin{tabular}{lrrrrrr}
\toprule
Run &
\shortstack{Receipt P50\\(ms)} &
\shortstack{Receipt P95\\(ms)} &
\shortstack{Dispatch P95\\(ms)} &
\shortstack{Scheduled-to-receipt P95\\(ms)} &
\shortstack{Send span\\(s)} &
\shortstack{Drain\\(s)} \\
\midrule
P-1 & 54.538 & 156.680 &  90.335 & 189.591 & 70.080 & 1.090 \\
P-2 & 54.007 & 147.549 & 245.315 & 326.828 & 70.235 & 0.955 \\
P-3 & 55.001 & 141.424 & 780.727 & 851.742 & 70.771 & 1.115 \\
\midrule
Median & 54.538 & 147.549 & 245.315 & 326.828 & 70.235 & 1.090 \\
\bottomrule
\end{tabular}
\end{table*}
\begin{figure*}[t]
  \centering
  \includegraphics[width=\textwidth]{mixed-load.pdf}
  \caption{Current-build P observations. Left: ordinary-receipt P95
    in the first seven ten-second windows; line and band show three-run
    medians and minima--maxima. Right: eight targets in P-1, with
    compensation and repayment preceding adaptation resumption.
    Public events are observations, not internal commit timestamps.}
  \label{fig:mixed-load}
  \label{fig:c3-resolution}
\end{figure*}
\subsection{Historical Reading: Equivalent Answers and Explicit Costs}
\label{sec:eval-reader}

We evaluate the exception-review workflow in
Section~\ref{sec:protocol-reader}: whether revised records preserve economic
answers and the original complete block identity, and what local reading,
storage, and maintenance they cost. The paused-adaptation runs above
separately test that economic closure proceeds before representation.

The timed historical-reader series retains build \texttt{721800c}; the same-consumer functional extension uses the current build. The reader is an honest local replica that has verified and executed
prefix $H$, querying original coordinates $(h,j,k)$ within one
application read transaction $V_H$. A loads original history and uses
the permanent decision index; B uses existing Canonical revisions.
Both consult the same permanent decisions and obligations, returning
identical input identity, TxID, coordinates, amount, historical
funding/debit, decision/height, and obligation state at $H$.
Authorization comes from successful execution binding exact revisions,
not a CH equation or matching BlockID. This is not remote-server
authentication.


\noindent\textbf{Body compatibility versus economic answers.}
An extended archival-consumer test compares A and B with M, a plain
serialized materialized funding row keyed by the same original identity
and coordinates. M caches immutable funding and decision fields and reads
current obligation status from $V_H$; it does not adapt any commitment.
All three return equal economic answers before and after real repayment.
B additionally verifies its authorized body against the actual saved
original commit. Rebuilding that changed body with initial part openings
preserves the header hash but changes the part-set header and is rejected
by the same commit check. After repayment, the authorized body and M's
cached rows remain unchanged while the obligation becomes Recovered.
Table~\ref{tab:c3-function} records these functional results. M is a
test-only cached view, not a third timed backend.

\begin{table}[t]
  \centering
  \caption{Functional comparison for local archival review. A: original
    body plus decision index; M: ordinary cached funding rows; B: C3.
    All use the same executed prefix and authorization records.
    The revised-body row is C3's representation contract.}
  \label{tab:c3-function}
  \small
  \begin{tabular}{@{}p{0.65\columnwidth}ccc@{}}
    \toprule
    Capability & A & M & B \\
    \midrule
    Equal economic answers at $H$ & Yes & Yes & Yes \\
    Retain original coordinates and ID & Yes & Yes & Yes \\
    Original body validates original commit & Yes & Yes & Yes \\
    Reserve-debit reference inside a revised body that validates the same commit & No & No & Yes \\
    Consult public authorization and state & Yes & Yes & Yes \\
    \bottomrule
  \end{tabular}
\end{table}

Fixtures use real signatures, successful ABCI execution, disk-backed
bbolt/GoLevelDB, and frozen BlockStores. A has its own original
BlockStore, without B's backup-layout cost. Point queries load a block
and decode the selected transaction; A's block query scans its range
index once and builds a timed map. Neither strategy receives a
precomputed decision map. Pages are warm, with no decoded-block or
result cache; \texttt{GOMAXPROCS=16}, \texttt{GOGC=100}.

Three fresh processes each test four $(n,m)$ fixtures, with
$n\in\{32,256\}$ payments and $m\in\{1,8\}$ compensated slots per block.
After 20 warm-ups per strategy, 300 interleaved queries per strategy
for compensated points, ordinary points, and whole blocks yield
21,600 timed queries. Timing includes read-transaction boundaries,
loading/decoding, record association, and result construction.
All timed reads use represented, materialized $H=4$; other stages
and subsequent real repayments are functional-equivalence checks,
not extra timing samples.

\begin{table*}[t]
  \caption{Historical reading and representation costs. A is original
    history plus decision index; B is Canonical. Point reads target
    compensated slots. Query entries are medians of three run-level
    P50s; cost entries are medians of three totals. Extra KV excludes
    common history, decisions, and application commit metadata.}
  \label{tab:eval-reader}
  \centering
  \small
  \setlength{\tabcolsep}{8pt}
  \begin{tabular}{@{}crrrrrr@{}}
    \toprule
    $(n,m)$ &
    \multicolumn{2}{c}{Point ($\mu$s)} &
    \multicolumn{2}{c}{Whole block ($\mu$s)} &
    \shortstack{Extra net KV\\(MiB/block)} &
    \shortstack{Local work\\(ms/block)} \\
    \cmidrule(lr){2-3}\cmidrule(lr){4-5}
    & A & B & A & B & & \\
    \midrule
    $(32,1)$  & 162.54  & 148.38 & 1,028.29 & 1,005.79 & 0.279 & 43.32 \\
    $(32,8)$  & 169.21  & 156.62 & 1,127.46 & 1,083.71 & 0.314 & 94.46 \\
    $(256,1)$ & 1,005.21 & 881.38 & 8,441.58 & 8,316.00 & 2.153 & 65.32 \\
    $(256,8)$ & 1,006.25 & 887.25 & 8,554.17 & 8,407.12 & 2.189 & 214.40 \\
    \bottomrule
  \end{tabular}
\end{table*}

Canonical reduces compensated-point P50 by 7.44--12.32\%
(12.6--123.8\,$\mu$s); whole-block differences are 1.49--3.88\%
(Table~\ref{tab:eval-reader}). In $(256,8)$ run 1, A performs 12
BlockStore Gets returning 622,183\,B; B performs none but reads
825,526\,B from application state. Mean per-query heap allocation is approximately
4.95 versus 2.24\,MB, while record-association medians are
11.50 versus 10.83\,$\mu$s. This indicates a layout/reconstruction
benefit, not fewer total bytes or cryptographic acceleration.
Decoded-object caching could change that tradeoff.

Representation adds 0.28--2.19\,MiB of net logical KV per repaired
block, for revisions, tasks, and backup/history records, net of deleted Pending. The $(256,8)$ 5,056-B batch is already
included. Local sequential adaptation, assembly, Finalize, synchronous
application Commit, and materialization total 43--214\,ms.
It excludes network/consensus waits, public-block saving, and
commit-signature construction. Common compensation costs are reported separately.
These figures characterize incremental repair on a CH-enabled build;
they do not isolate the steady-state cost of CH encoding against a
non-CH build.



Equivalence checks cover compensation before representation,
representation before materialization, and repayment preserving the
historical debit with state Recovered. Invalid coordinates,
cross-prefix decisions, and wrong debit/amount are rejected;
undecodable revisions fail rather than fall back. These checks do
not establish detection of arbitrary decodable corruption.
Canonical offers authorized funding representation at stable
historical coordinates, with measured maintenance costs, rather
than an economically necessary or universally faster alternative
to indexed history.

\subsection{Injected Latency and Single-Fault Runs}
\label{sec:eval-faults}
On \texttt{9879f64}, four conditions run three times each in alternating
order. H uses the injection proxies without added delay; L adds 25\,ms
per direction to member RPC and committee P2P traffic. L-M and L-C
add that delay and suspend member~3 or committee validator~3,
respectively, from 15 to 45\,s. Each run sends 100 warm-up payments,
1,200 ordinary payments at 20/s, and two 10-hop chains starting inside
the fault window. The actual health-RPC RTT rises from 1.17 to
52.20\,ms; both RPC and P2P mean holds are about 25.2\,ms per direction.
The supplement specifies the FIFO injection, fixed resources and per-run results.

All 15,840 payments close, with equal committee states, empty outboxes
and zero public gap. No compensation is needed.
Table~\ref{tab:eval-faults} separates receipt, public observation and
observation of all four members. With one member suspended, new QCs
use the other three members and public completion continues, while
the absent member's observation creates a peak of 630--641 outstanding
tasks. With a validator suspended, receipt P95 remains 108--111\,ms
and 400--510 payments reach public observation inside each fault
window, but public P95 rises to 9.5--12.6\,s. Retained commit
certificates verify five paused-validator round-0 proposer opportunities
per run; each commits in round~1 without that validator's signature.
Both chains finish certification within each fault window, and every
run drains within 1.25\,s of its last ordinary send. These are
single-host injected-delay results, separate from the fast/wait comparison.

\begin{table}[t]
\caption{Injected latency and single faults: medians of three run-level
statistics over all 1,200 ordinary payments per run; warm-up and chains are excluded. Public and member columns are send-to-observation P95;
the member column waits for all four members.}
\label{tab:eval-faults}
\centering\small
\begin{tabular}{@{}lrrr@{}}
\toprule
Condition & Receipt P50/P95 (ms) & Public (s) & Members (s)\\
\midrule
H   & 41.56/83.20  & 1.28 & 1.36\\
L   & 81.91/118.22 & 1.32 & 1.41\\
L-M & 77.71/112.85 & 1.30 & 29.00\\
L-C & 77.41/109.32 & 10.02 & 10.04\\
\bottomrule
\end{tabular}
\end{table}


\subsection{Boundary and Composition Checks}
\label{sec:eval-boundaries}
Deterministic tests on \texttt{9879f64} use actual member approvals,
ABCI execution, and authenticated follower updates. They reproduce
no-QC capacity fragmentation, a same-input 2/2 split, and two
Intent-conflict compensation rounds with complete CAL/FUEL accounting.
An external top-up resolves the first case but not the second. Real
threshold adaptation covers both representation/repayment orders;
retained-state restart reproduces each original command history.
A commit-before-signature interruption preserves locks and reservations
and permits an idempotent retry. The supplement gives the per-prefix
ledgers, negative-entry tests, and restart assertions. The current-build
proof-to-code map identifies the entry points, atomic guards, and existing
regressions supporting each property. An additional raw-byte ABCI test
rejects unrelated or forged QCs and forged owner signatures at CheckTx,
ProcessProposal, and FinalizeBlock without business-state changes.

\noindent\textbf{Archived evidence.}
The supplement retains earlier continuation, throughput, capital,
organization-path, and availability experiments under their original
builds and configurations. They provide historical context, not
measurements of \texttt{f856c7b}; cross-version differences are not
attributed to individual changes.
